%% kapa: exam=1 round=21 year=2010 period=2 form=A subject=appraisal_law subject_ko=부동산관계법규 seq=28 number=68
%% source: Q-Net 감정평가사 기출문제 (cst00309) artlSeq=5024181 file=extracted/1008/5024181/2010년 제21회 감정평가사 1차시험 기출문제/2교시.hwp
%% conversion: kordoc 4.24.51 -> tools/gen1.py (auto); review=proofread-gpt-sol; flags=none
\begin{question}{68}
국유재산법령상 일반재산의 매각․교환 및 양여에 관한 설명으로 옳은 것은?
\begin{choices}
  \item 일반재산을 용도를 지정하여 매각하는 경우에는 그 재산의 매각일로부터 10년 이상 지정된 용도로 활용하여야 한다.
  \item 일반재산의 매각대금은 분할 납부할 수 없다.
  \item 일반재산의 매각으로 인한 소유권이전은 매각대금의 완납 이전에는 할 수 없다.
  \item 관리청등이 일반재산을 직접 공용이나 공공용으로 사용하려는 지방자치단체에게 양여하는 경우에는 총괄청과 협의할 필요는 없다.
  \item 관리청등은 일반재산을 교환하려면 그 내용을 기획재정부장관에게 보고하여야 한다.
\end{choices}
\end{question}
\answer{1}
